%%
%% This is file `sample-manuscript.tex',
%% generated with the docstrip utility.
%%
%% The original source files were:
%%
%% samples.dtx  (with options: `manuscript')
%% 
%% IMPORTANT NOTICE:
%% 
%% For the copyright see the source file.
%% 
%% Any modified versions of this file must be renamed
%% with new filenames distinct from sample-manuscript.tex.
%% 
%% For distribution of the original source see the terms
%% for copying and modification in the file samples.dtx.
%% 
%% This generated file may be distributed as long as the
%% original source files, as listed above, are part of the
%% same distribution. (The sources need not necessarily be
%% in the same archive or directory.)
%%
%% Commands for TeXCount
%TC:macro \cite [option:text,text]
%TC:macro \citep [option:text,text]
%TC:macro \citet [option:text,text]
%TC:envir table 0 1
%TC:envir table* 0 1
%TC:envir tabular [ignore] word
%TC:envir displaymath 0 word
%TC:envir math 0 word
%TC:envir comment 0 0
%%
%%
%% The first command in your LaTeX source must be the \documentclass command.
%%%% Small single column format, used for CIE, CSUR, DTRAP, JACM, JDIQ, JEA, JERIC, JETC, PACMCGIT, TAAS, TACCESS, TACO, TALG, TALLIP (formerly TALIP), TCPS, TDSCI, TEAC, TECS, TELO, THRI, TIIS, TIOT, TISSEC, TIST, TKDD, TMIS, TOCE, TOCHI, TOCL, TOCS, TOCT, TODAES, TODS, TOIS, TOIT, TOMACS, TOMM (formerly TOMCCAP), TOMPECS, TOMS, TOPC, TOPLAS, TOPS, TOS, TOSEM, TOSN, TQC, TRETS, TSAS, TSC, TSLP, TWEB.
% \documentclass[acmsmall]{acmart}

%%%% Large single column format, used for IMWUT, JOCCH, PACMPL, POMACS, TAP, PACMHCI
% \documentclass[acmlarge,screen]{acmart}

%%%% Large double column format, used for TOG

\documentclass[acmtog, authorversion, tikz, border=10pt]{acmart}
\usepackage[htt]{hyphenat}
\usepackage{graphicx}
\graphicspath{{images/}}
\usepackage{pgfgantt}
\usepackage{tikz}
\usetikzlibrary{arrows.meta, bending, positioning, calc}
\usepackage{algorithm}
\usepackage{algpseudocode}
\usepackage{todonotes}
\usepackage{pgfplots}
\pgfplotsset{compat=1.18}
\usepgfplotslibrary{fillbetween}
\usepackage{amsmath}


% continuation indent patch http://tex.stackexchange.com/questions/78776/forced-indentation-in-algorithmicx
\newlength{\continueindent}
\setlength{\continueindent}{2em}
\usepackage{etoolbox}
\makeatletter
\newcommand*{\ALG@customparshape}{\parshape 2 \leftmargin \linewidth \dimexpr\ALG@tlm+\continueindent\relax \dimexpr\linewidth+\leftmargin-\ALG@tlm-\continueindent\relax}
\apptocmd{\ALG@beginblock}{\ALG@customparshape}{}{\errmessage{failed to patch}}
\makeatother
% end continuation indent patch


%%%% Generic manuscript mode, required for submission
%%%% and peer review
%\documentclass[manuscript,screen,review]{acmart}
%% Fonts used in the template cannot be substituted; margin 
%% adjustments are not allowed.
%%
%% \BibTeX command to typeset BibTeX logo in the docs
\AtBeginDocument{%
  \providecommand\BibTeX{{%
    \normalfont B\kern-0.5em{\scshape i\kern-0.25em b}\kern-0.8em\TeX}}}

%% Rights management information.  This information is sent to you
%% when you complete the rights form.  These commands have SAMPLE
%% values in them; it is your responsibility as an author to replace
%% the commands and values with those provided to you when you
%% complete the rights form.

\setcopyright{acmcopyright}
\copyrightyear{2022}
\acmYear{2022}

%% These commands are for a PROCEEDINGS abstract or paper.
\acmConference[TScIT 45]{45$^{th}$ Twente Student Conference on IT}{July 3,
  2026}{Enschede, The Netherlands}
%
%  Uncomment \acmBooktitle if th title of the proceedings is different
%  from ``Proceedings of ...''!
%
%\acmBooktitle{Woodstock '18: ACM Symposium on Neural Gaze Detection,
% June 03--05, 2018, Woodstock, NY} 
%\acmPrice{15.00}
%\acmISBN{978-1-4503-XXXX-X/18/06}


%%
%% Submission ID.
%% Use this when submitting an article to a sponsored event. You'll
%% receive a unique submission ID from the organizers
%% of the event, and this ID should be used as the parameter to this command.
%%\acmSubmissionID{123-A56-BU3}

%%
%% For managing citations, it is recommended to use bibliography
%% files in BibTeX format.
%%
%% You can then either use BibTeX with the ACM-Reference-Format style,
%% or BibLaTeX with the acmnumeric or acmauthoryear sytles, that include
%% support for advanced citation of software artefact from the
%% biblatex-software package, also separately available on CTAN.
%%
%% Look at the sample-*-biblatex.tex files for templates showcasing
%% the biblatex styles.
%%

%%
%% The majority of ACM publications use numbered citations and
%% references.  The command \citestyle{authoryear} switches to the
%% "author year" style.
%%
%% If you are preparing content for an event
%% sponsored by ACM SIGGRAPH, you must use the "author year" style of
%% citations and references.
%% Uncommenting
%% the next command will enable that style.
%%\citestyle{acmauthoryear}

%%
%% end of the preamble, start of the body of the document source.
\begin{document}

%%
%% The "title" command has an optional parameter,
%% allowing the author to define a "short title" to be used in page headers.
\title{DeltaParq -  Using Delta Compression for Efficient Storage of OpenINTELs Historical DNS Records}

%%
%% The "author" command and its associated commands are used to define
%% the authors and their affiliations.
%% Of note is the shared affiliation of the first two authors, and the
%% "authornote" and "authornotemark" commands
%% used to denote shared contribution to the research.

\author{Tristan Bottenberg}
\email{t.d.bottenberg@student.utwente.nl}
\affiliation{%
  \institution{University of Twente}
  \streetaddress{P.O. Box 217}
  \city{Enschede}
  \country{The Netherlands}
  \postcode{7500AE}
}



%%
%% By default, the full list of authors will be used in the page
%% headers. Often, this list is too long, and will overlap
%% other information printed in the page headers. This command allows
%% the author to define a more concise list
%% of authors' names for this purpose.
%% \renewcommand{\shortauthors}{Authors}

%%
%% The abstract is a short summary of the work to be presented in the
%% article.
\begin{abstract}
    The OpenINTEL project gathers millions of records from the Domain Name System (DNS) 
    for the purpose of historical DNS research. This research shows that a big part of 
    the DNS remains unchanged and based on this presents a delta-based storage system, DeltaParq. 
    The DeltaParq system only stores the changes compared to a snapshot instead of
    all daily retrieved data points to improve on storage efficiency. The DeltaParq is then 
    evaluated to show the effects of different delta-based compression strategies on both 
    storage and analysis efficiency.
\end{abstract}

%%
%% Keywords. The author(s) should pick words that accurately describe
%% the work being presented. Separate the keywords with commas.
\keywords{DNS, Delta-Compression, Big Data}

\settopmatter{printacmref=false}

%% A "teaser" image appears between the author and affiliation
%% information and the body of the document, and typically spans the
%% page.
%% \begin{teaserfigure}
%%   \includegraphics[width=\textwidth]{sampleteaser}
%%   \caption{Seattle Mariners at Spring Training, 2010.}
%%   \Description{Enjoying the baseball game from the third-base
%%   seats. Ichiro Suzuki preparing to bat.}
%%   \label{fig:teaser}
%% \end{teaserfigure}

%%
%% This command processes the author and affiliation and title
%% information and builds the first part of the formatted document.
\maketitle

\section{Introduction}

Since its introduction in 1983 as a small name system on the ARPANET, 
the Domain Name System (DNS) has grown to 392.5 million registrations 
across all top-level domains (TLDs)~\cite{dnib_q1_2026}. 
The OpenINTEL project\footnote{\url{https://openintel.nl/}} actively records daily historical data 
of most TLDs. This historical data has been essential for research in many domains 
such as cybersecurity~\cite{Sommese_2024} and infrastructure resilience~\cite{10.1145/3517745.3561458}.
The original design of the OpenINTEL infrastructure~\cite{7460220} stores all data in the 
Parquet\footnote{\url{https://parquet.apache.org/}} columnar storage format for each day and TLD. \\

This research is motivated by the hypothesis that a majority of the collected DNS records 
remain unchanged between measurements. Under this hypothesis a majority of the data stored 
every day by the OpenINTEL project is a duplicate of the previous day. Instead of storing 
all data every day, using a delta-based storage solution, which only stores the records 
that changed compared to a previous snapshot, could substantially reduce the storage usage. 
Furthermore, as the size of the stored records decreases, less data needs to be retrieved 
from storage which could improve the analysis efficiency. The downside of using delta-based 
storage however, is that reconstructing a data point requires reconstructing the full 
dependency chain, which may decrease the analysis efficiency.

Designing such a delta-based storage system presents several non-trivial challenges.\\
Since the data is stored in a format where row-group coherence matters, where it is important 
which row and column a value belongs to, it is not possible to store just the new values of each 
column as that results in not knowing which rows these values might belong to. 
This constrains the system to copy either the full column or the full row into the delta. 
A full-column delta however is impractical for this dataset: a single row receiving a new value would 
require storing all millions of rows in the delta column, yielding no storage benefit. A full-row 
delta follows as the natural choice, but this introduces further problems, namely volatile columns.
The OpenINTEL database contains several volatile columns that change almost every 
measurement, such as the measurement timestamp and round-trip-time. In a row-delta implementation this 
would cause almost all rows to be copied in full every day resulting in no improvement in 
data storage efficiency. However, in the case of a column-delta implementation, a singular row 
receiving a new value would result in all millions of rows being stored in the delta, again resulting 
in almost no storage benefit. From this follows that the system needs a decomposition of the database 
schema to be able to handle both the volatile data while also being able to use delta-compression 
for the non-volatile data.\\

Finally, multiple strategies can be used for a delta-based storage system. For example, 
progressive deltas store deltas based on the previous measurement while snapshot deltas 
would store all deltas based on a single snapshot. These different strategies present tradeoffs
between storage efficiency and reconstruction cost. The exact causes of these tradeoffs lie in 
the characteristics of the underlying data of the OPENIntel system. \\ 

The main focus of this research is the impact of using a delta-based storage system, and its different 
strategy designs on storage and analysis efficiency. This leads to the following research question:

\begin{center}
    \textit{\textbf{
    To what extent can storage and analysis efficiency be improved by transforming a daily historical DNS 
record database into a delta-based storage system?}}
\end{center}

To answer the research question, this paper will first investigate the state of the art in delta-based systems 
to identify the gaps that motivate this work. Furthermore, the paper will evaluate the correctness of the underlying 
hypothesis. Based on the explored system constrains and research gaps, DeltaParq is introduced, a schema-aware delta-based 
storage prototype built for the OPENIntel project. An evaluation framework will be presented that assesses both storage 
and analysis efficiency. With this framework multiple delta strategies will be assessed of which the results will be presented.
Finally, using the results, conclusions are drawn on the practical viability of a delta-based storage system for a daily DNS 
record database.


\section{State of the Art}

\subsection{Delta Storage Systems}

Several programs have been developed for adding versioning and incremental updates 
to the write-only Parquet storage files. 

The program closest to meeting the needs of the OpenINTEL storage system is 
Apache Hudi\footnote{\url{https://hudi.apache.org/}}. Hudi is a storage management 
system built on top of database formats like Parquet. Hudi adds 
log files onto these files, which describe changes and can be viewed 
as deltas to the original Parquet file. However, it is limited by its retrieval and 
storage capabilities. Hudi was created as an efficient way to apply edits to a 
write-only database. Because of this, it only has functions to retrieve the current status 
which applies all delta logs at once. Hudi lacks the feature of constructing the table for 
any given snapshot. Furthermore, a built-in feature that cannot be removed from Hudi is its 
Compaction feature. Hudi will regularly update the base snapshot which merges the first few 
log files into the base and writes a new base file. This would destroy the history stored in 
the log files.

Both Apache Iceberg\footnote{\url{https://iceberg.apache.org/}} and Delta Lake\footnote{\url{https://delta.io/}} 
take a different approach to extending the Parquet storage format. Both systems create a time-travel 
system by keeping track of snapshots for every update. This is done by creating new Parquet 
files for every update to the base. A timestamp can then be added to the query to which these systems 
retrieve the relevant Parquet files. However, the important pitfall with these systems is that 
they are not capable of storing deltas; they only store snapshots. Unlike Hudi, these systems 
are not designed to progressively add snapshots together, but for full snapshots. These systems 
expect every measurement to be a full snapshot repeating all the data and not saving on storage 
efficiency at all. If these systems would be used by adding only the changed rows to each day, 
the systems are not capable of combining the different snapshots together, this would have to be 
done manually where each Parquet for each day would be retrieved and combined.

Furthermore, none of these current systems are able to handle formats with partially 
volatile columns that update every measurement.

\subsection{DNS Measurement Storage Systems}

Next to the OpenINTEL project which this research will build on, several DNS 
history providers exist. The most prominent one is the ActiveDNS 
project\footnote{\url{https://activednsproject.org/}}. Similar to OpenINTEL, ActiveDNS stores 
daily DNS records. The ActiveDNS project collects data from multiple vantage points; 
this results in possible duplicate data which is deduplicated in the storage module 
of the system~\cite{DBLP:conf/raid/KountourasKLCND16}. However, like OpenINTEL, 
ActiveDNS does not have any deduplication of data between separate days. 

Furthermore, commercial providers exist such as the Farsight 
DNSDB\footnote{\url{https://www.domaintools.com/}}, WhoisXMLAPI's DNS 
Database\footnote{\url{https://dns-history.whoisxmlapi.com/database}}, and SpamHaus's
passive DNS database\footnote{\url{https://www.spamhaus.com/data-access/passive-dns-api/}}. 
None of these providers provide any public documentation on their storage system, thus 
it cannot be concluded if any type of delta-compression is used to efficiently store their data.

\subsection{Delta Strategies}

Delta storage is used in many video codecs used to efficiently store video data, such as the MPEG video encoding 
standards~\cite{isoiec11172}. The MPEG codecs store frames of a video as a combination 
of base images, forward predicted images, and bi-directionally predicted images. The 
bi-directionally predicted images allow for even smaller deltas as the similarities to both 
the previous and the next data point can be used. Under the hypothesis that DNS records 
do not change often, this could be beneficial for the OpenINTEL storage system in principle.
However, the OpenINTEL storage system does not have immediate access to all data. 
The system would not be able to store the current day's data without knowing the data of the 
next day. This would cause more temporary storage overhead and adds complexity to the 
storage system which is unfavorable.

Furthermore, Molfetas et al.~\cite{DBLP:conf/awc/MolfetasWZ14} tested several combinations of files and strategies 
with the goal of reducing the amount of dependencies on other snapshots while retaining storage efficiency. 
This work is similar in nature since it touches on the effects of using delta-compression, however it does not take into account 
the vast amount of the data being stored by the OpenINTEL project and the significant time loss of 
I/O operations on the hard disks for retrieval of the resulting data that this creates.

LoopDelta, a delta-based compression framework with backups in mind, does identify the extra I/O 
operations necessary for reconstruction of delta files as a bottleneck in delta-based 
storage~\cite{10.1145/3721485}. But, similar to the aforementioned video codecs, LoopDelta relies 
partially on backwards deltas which require the next day to be known in the OpenINTEL system 
for the compression optimization.


\section{Evaluating The Hypothesis}

\begin{table}
    \caption{Total changing rows in a month in the $.li$ and $.ch$ datasets.}
    \label{tab:measurement_differences}
    \begin{center}
        \begin{tabular}[c]{|l|r|r|r|} \hline
            Dataset & Changed Rows & Total Rows & Percentage \\ \hline
            .li & 5182691 & 32324410 & 16.0\% \\
            .ch & 280111056 & 1570870201 & 17.8\% \\
            
            \hline
        \end{tabular}
    \end{center}
\end{table}

To validate the hypothesis set out in the introduction that a majority 
of the DNS records remain unchanged between measurements, a set-based 
comparison can be defined. Let $D_n$ be the complete set of measurements 
on day $n$. The set of changed records, which are all records present in 
$D_n$ and not present in $D_{n-1}$ can be denoted as:
$$\Delta_n = D_{n} \setminus D_{n-1}$$
Table \ref{tab:measurement_differences} shows the results of this set 
operation over the span of May 2026 on both the Liechtenstein and the Switzerland dataset. 
These results support the hypothesis, on average only about 17\% of the rows change per measurement.



\section{DeltaParq System Design}

\subsection{System Overview}

\subsubsection{Infrastructure}
For the design of the DeltaParq system, the structure and resources 
of the current system has been taken into account. The OPENIntel system 
utilizes a Hadoop file system (HDFS) running on top of an S3 storage cluster\cite{7460220}.
Furthermore, Apache Spark, with its Python derivative psypark, is used for 
querying and analyzing the dataset. While the DeltaParq system has been 
designed for the OPENIntel project, it has also been designed with 
generalizability in mind. The system uses a configuration file to set 
which columns are part of the primary key(PK) and which columns are volatile, 
this way other datasets could utilize the DeltaParq as well. All join operations 
within the system utilize this PK as the join condition. Furthermore, 
the integrations of the DeltaParq are extendible. Integrations with different 
file systems are easy to integrate and can after integration easily be switched 
in the configuration.


\subsubsection{Self-Describing Storage}

Like the current OPENIntel system, DeltaParq stores data using a Hive-style path on an S3 
bucket. DeltaParq constructs the path using the TLD, date, and S3 bucket that has been set 
in the configuration.
$$"s3a://{bucket}/tld={TLD}/year={year:04}/month={month:02}/day={day:02}/"$$
\todo[inline]{Fix this formatting}
Using this hierarchy, DeltaParq treats each month as a self-contained unit, this way each delta 
can only depend on a day in the same month. As the first day of the month is always the start, a 
part of the configuration called the $month\_config$ is stored in the month's directory. This part of  
the config describes the primary key, volatile columns, and strategy. This makes each stored month 
self-describing and ensures each dataframe can be correctly reconstructed without relying on the 
current system's configuration.

\subsubsection{Storage Strategies}

Since a delta will always be based on a single base dataframe, a storage strategy which determines 
which delta is based on which base, a strategy can be encoded as a function $strategy: \mathbb{Z} \rightarrow 
\mathbb{Z}$. Furthermore, a day is stored as a full snapshot if $strategy(d) = d$.

\subsection{Putting Data into the DeltaParq}

\subsubsection{Schema Decomposition}

As mentioned in the problem statement, the OPENIntel dataset 
contains both volatile and non-volatile data. As the volatile data is almost guaranteed 
to change between each measurement, these columns need to be removed from the non-volatile 
data using a schema decomposition. Using pyspark, the decomposition can be easily done by 
selecting only the necessary columns of a dataframe. For the volatile data, both the columns 
marked as volatile in the configuration and the columns of the primary key are selected 
such that this data can be joined with the non-volatile data for which all original columns, 
without the columns marked as volatile, but including the primary key, are selected. 
The resulting volatile dataframe can immediately be written to storage, while non-volatile 
dataframe needs to be processed using delta generation to be stored in the DeltaParq.

\subsubsection{Delta Generation}

The concept of delta generation is based on filtering out all the rows that are different 
in the updated version compared to the base version. The OPENIntel dataset, and DNS records 
in general, have an inherent flaw making this difficult: DNS records do not have a unique identifier.
The identifiers for a DNS record is its domain name and record type, but this combination does not have to be
unique. For example, the domain name example.com could return multiple TXT records, if between measurements
one TXT record changes while another is removed, there is no way of keeping track which of the TXT records 
changed and which one got removed. Not having a unique identifier leaves two options. Either every change is tracked 
as an addition or subtraction, this would result in two subtractions and one addition being stored in the above example, 
or the whole set of rows identified by the domain name and type is deemed as updated and thus stored in the delta.
For the DeltaParq, the latter option of storing the full identified block was chosen. \todo[inline]{explain why!!!}
A simple subtraction based algorithm, such as Algorithm~\ref{alg:subtract-based-delta} uses the pyspark subtract function 
to gather all rows in the updated dataframe not present in the removed dataframe. Rows that do not exist in the update 
but do exist in the base dataframe will not show up in this subtraction, thus a second subtraction needs to be done 
to gather all rows in the base dataframe not present in the update. Combining these keys together 
then results in all keys necessary for the delta.

\begin{algorithm}
    \caption{Subtract-based Delta Generation}
    \label{alg:subtract-based-delta}
    \begin{algorithmic}[1]
        \Function{SubtractDelta}{$base$, $update$, $PK$}
            \State $\Delta_{new} \gets \pi_{PK}(update \setminus base)$
            \State $\Delta_{removed} \gets \pi_{PK}(base \setminus update) \bowtie \pi_{PK}(update)$
            \State $\Delta_{keys} \gets \Delta_{new} \cup \Delta_{removed}$
            \State \Return $update \bowtie \Delta_{keys}$
        \EndFunction
    \end{algorithmic}
\end{algorithm}

While the subtract-based delta generation worked, it is not efficient. The main bottleneck of the 
OPENIntel system is the retrieval from storage. If both the dataframes do not fit in the allocated RAM,
during operations, they need to be retrieved multiple times for the subtractions after which the update 
dataframe is needed once more for the join with the delta keys. To counter these retrievals, a hash-based 
delta generation algorithm, Algorithm~\ref{alg:hash-based-delta}, has been designed. The hash-based delta
generation uses an auxiliary function to transform a dataframe such that each row is a primary key combined 
with a hash generated from all values in all rows from a block. This way the data stored in memory is smaller 
and only requires one match between dataframes instead of comparing each row with all others like the 
subtract-based generation would. Using this auxiliary function a hash-based dataframe can be generated 
for both the base and the update to then join these dataframes and filter out all keys that have matching 
hashes. Finally, the primary keys that remain will identify all rows from the update to store in the delta.

\begin{algorithm}
    \caption{Hash-based Delta Generation}
    \label{alg:hash-based-delta}
    \begin{algorithmic}[1]

        \Function{HashBlock}{$dataframe$, $PK$}
            \State $dataframe' \gets dataframe$ \textbf{with} $row\_hash \gets H(row)$ \textbf{for each} $row$
            \State \Return $\text{groupBy}(dataframe', PK,\ H(\text{sort}(row\_hashes)))$
        \EndFunction
        \State
        \Function{GenerateHashDelta}{$base$, $update$, $PK$}
            \State $base\_hashes \gets \text{HashBlock}(base, PK)$
            \State $update\_hashes \gets \text{HashBlock}(update, PK)$

            \State $joined \gets update\_hashes \overset{left}{\bowtie} base\_hashes$

            \State $\Delta_{keys} \gets \sigma_{\ base.hash = \emptyset\ \vee\ base.hash \neq update.hash}\ (joined)$

            \State \Return $update \bowtie \Delta_{keys}$
        \EndFunction

    \end{algorithmic}
\end{algorithm}

\subsection{Retrieving Data from the DeltaParq}

Retrieving a dataframe for a given day from the DeltaParq requires exploring the full 
dependency chain, reconstruct all dataframes in the dependency chain, and finally join 
the dataframe with the volatile data of the day that is to be retrieved.

\subsubsection{Dependency Resolution}

Given a delta strategy function $strategy: \mathbb{Z} \rightarrow \mathbb{Z}$ that maps each
day to the day it depends on, the dependency chain for a given day can be resolved for a day by 
walking the chain until a snapshot, identified by $strategy(d) = d$, is reached. If multiple 
days are requested, the dependency chain is walked for each day. If at any point a dependency is 
identified that was already part of the chain of a different day, the chains are merged. This way 
the dependency resolver always returns a chain without duplicates that ensures that each day's 
dependency precedes itself in the chain:
$$\forall d \in chain,\quad strategy(d) = d \;\lor\; strategy(d) \in chain_{[:i_d]}$$
where $chain_{[:i_d]}$ denotes the chain up until $d$.

\subsubsection{Delta Reconstruction}

Using the resolved dependency chain, the full dataframe for each day in the chain can be reconstructed 
which will result in all the requested days being reconstructed as well. For a snapshot day the data is read 
directly from storage, while for a delta day the delta is applied to a previously reconstructed dataframe. 
The application of a delta to a base requires filtering out all primary keys in the delta from the base as 
these have changed between the two measurements using an anti join. Afterward, the result can be unioned with the delta to 
construct the full dataframe:
$$(base \triangleright delta) \cup delta$$
As a final step for reconstructing the full dataframe, the applied delta needs to be joined with the volatile dataframe. 
Even if the volatile columns are not used in the further analysis, the volatile dataframe determines which data was present 
for the specific measurement. Since the delta only stores updates found in the new measurement, blocks that were fully removed 
from in the update cannot be present in the delta dataframe. Since the volatile dataframe measures all the present blocks and 
its volatile columns for each measurement, this can be used to filter out all the completely removed blocks that are still 
present in the applied delta as these would not be present in the volatile dataframe.

\subsubsection{Public Interface}

The DeltaParq offers both a $get()$ and $get\_list()$ interface for retrieving data. With $get\_list()$ the retrieval 
of the dataframes is structured such that any dataframes that would be retrieved multiple times in the different dependency 
chain of each requested day is only retrieved once. Furthermore, the retrieval methods of the DeltaParq offer a column selection 
and filter pushdown. Both the column selection and filter are immediately applied to the dataframes retrieved from storage, like 
pyspark's lazy evaluation would do. This way only the necessary data is retrieved speeding up the retrieval and analysis process.

\subsection{Performance Constraints}

The DeltaParq would benefit significantly from data being physically partitioned by the primary key on write. Both the store 
and retrieval part of the DeltaParq use various joins and substracts which are wide operations. In the case of spark executing 
on multiple excutors, the executors need to share information for these wide operations. An executor may have a row that needs to be joined with a row 
from the other dataframe present on a different executor causing waiting time on other nodes and extra overhead for sharing the data. 
If all data is partitioned such that each partition is guaranteed to have all rows for each primary key it has, no data sharing is 
necessary for these wide operations as each executor now knows they aren't possibly missing data. If this partitioning is done efficiently
on write and pyspark were able to read this, the now broad execution path of combining all the data can be transformed into a direct linear 
path for each executor.

However, while pyspark does support saving partitions in buckets, it is only able to do this by using a Hive metastore. The 
metastore would store all the information on which buckets were used to ensure correctness, the same way DeltaParq uses a stored 
monthly configuration for its self-describing storage. Without a metastore pyspark does not have a guarantee that the data is 
partitioned or bucketed and will always assume no partitioning. Furthermore, pyspark does not offer a way of telling the program 
the data is partitioned in a certain way without the metastore. In the current deployment, no metastore is available, therefore 
the DeltaParq has not been integrated with any form of partitions.


% Keynotes:


% build on top of spark/current used products
% schema decomposition
    % Use config to get the columns and just split that shit open
% delta generation
    % subtraction method
    % "faster" hash method
% delta joining
    % generate dependencies
    % sort dependencies such that base is always build before being used
    % return
% public methods
    % put / put_list
    % get / get_range (explain with/without explicit removal save and select/filter pushdown)
    % rebuild
    
% optimization problems with spark
    % no metastore == no partitioning/order
    % big sad


\section{Evaluation Framework}

To answer the main research question on the efficiency gains in both storage and analysis, an 
evaluation framework is presented that has the ability to measure the impact of using different 
versions of the DeltaParq system compared to a base version of the system as it is currently 
used.

\subsection{Measurement Metrics}

The primary metrics measured are the total storage size for storing a month worth of data, which 
was determined as the maximum self-contained unit in the system design, and the analysis time for each 
day of the month given a specific analysis method. The storage size for the full month will be able to,
once compared to the storage size of a base version, determine the storage efficiency of each version, 
while the analysis time for each day of the month will be able to determine, compared to the same base 
version, how much the analysis efficiency has increased or decreased. 

While these two metrics capture the full scope necessary for answering the research question. The research 
is held to determine if a delta-based storage system is a viable option for the OPENIntel dataset. From this 
follow two extra metrics. First the time to construct the delta-compressed version which will be 
used to determine the effect of using delta-compression on the daily task of gathering and storing all the new 
data. Secondly, the time to reconstruct the full original data will be measurement to report the effects that may 
occur when needing the full data for possible migration or other use cases.

% The client commisioning the research, the OPENIntel project, 

\subsection{Measurement Validity}

\subsection{Measurement Methodology}


\section{Experiments \& Results}


\subsection{Strategies}


\section{Conclusions and Future Work}

\subsection{Conclusions}
This research investigated to which extent the storage and analysis efficiency can be improved by transforming a daily
historical DNS record database into a delta-based storage system. The results confirmed the underlying hypothesis that 
a majority of the collected DNS records remain unchanged between measurements. The introduced DeltaParq, a schema-aware 
delta-based storage prototype build for the OpenINTEL project, showed that this can be used to reduce storage usage.
An evaluation over multiple datasets showed that up to about 49.25\% of storage space can be saved using the 
Progressive strategy, which also comes with a significantly longer analysis time with an increase of 98.65\%. Other evaluated strategies, 
such as the BST strategy, show that saving a substantial amount of storage space can be done with less impact 
on the analysis time. Using this the system can be tuned to the preferences of its user.

% Hypothesis deemed correct
% Succesfully constructed the DeltaParq prototype
% Maximum gain lies around 50% but comes with a loss
% However storage gain definetly possible with less removal from analysis time (BST)

\subsection{Future Work}

The findings of this research point to multiple directions for future research. Firstly, the bottleneck identified in Section~\ref{sec:performance_constraints}, 
the shuffle overhead caused by PySpark not being able to use sorted or partitioned data without a Hive metastore. A deployment with support for partitioning on storing and reading may significantly 
reduce the overhead in analysis time, which could potentially make the analysis time of strategies that use long dependency chains, such as 
the Progressive strategy, significantly lower. Furthermore, it may make more frequent anchor resets, such as weekly instead of monthly, 
more viable as the increased analysis time identified in~\ref{sec:discussion} would have less of an impact.

Secondly, this research evaluated the DeltaParq prototype on only 2 datasets based on a single underlying structure. Evaluating the 
schema-aware delta-based storage systems over more datasets, both within the same OpenINTEL structure, and other domains with similar 
volatile and non-volatile data columns, would strengthen the conclusions and the generalizability beyond the OpenINTEL system.

Finally, strategies can be optimized for both better performance and data safety.
Further research could explore the possibility of using the structure of the underlying dataset combined with the 
properties of the strategies to automatically select a strategy that performs optimally for a user defined balance 
between storage reduction and analysis time increase. Furthermore, certain strategies could be extended. For example,
adding a bridging delta between the last day of the month and the first day of the next month in the progressive strategy, 
would allow for a smart system to recover a singular corrupted delta.


% Spark optimization -> May make more aggressive strategies possible
% Test on more datasets
% Create a model dataset -> Strategy


%%
%% The next two lines define the bibliography style to be used, and
%% the bibliography file.
\bibliographystyle{ACM-Reference-Format}

\bibliography{paper.bib}

\end{document}
\endinput
%%
%% End of file `sample-authordraft.tex'.
